\section{Application to graph containment}\label{sec:graphs}

It remains to construct the tree required by Theorem~\ref{thm:tree}.
We use the normalized-containment maximization and conditional-extension
argument of Tran~\cite[Section~3, Lemma~3.1]{Tran26}, with constants suited
to the factor-$16$ separation of capacities. The copy-count identity below
also appears in Mossel, Niles-Weed, Sun, and Zadik~\cite[Equation~(5)]{MNSZ22}.

Put $q=\pE(n,H)$ and $\sigma=128q$. Delete the isolated vertices of $H$
and of each graph subsequently considered, and identify these graphs
with their edge sets. This does not change the event of containing
$H$: any copy of its nonisolated part can be extended by choosing the
required number of additional vertices, since $n\ge|V(H)|$ and
copies need not be induced. Every edge subset used below, on its
incident vertices, is a subgraph of the original $H$, so its
expectation at $q$ is at least $1/2$.
Throughout this deletion, $q$ remains the threshold defined using the
original graph $H$.

For a graph $F$ without isolated vertices, let
\[
 M(F)=N(K_n,F),\qquad M(I,F)=N(I,F).
\]
In particular, $M(F)=(n)_{v(F)}/|\operatorname{Aut}(F)|$, where
$(n)_v=n(n-1)\cdots(n-v+1)$ and $\operatorname{Aut}(F)$ is the group of
vertex permutations preserving adjacency. Dividing by this group counts
each ordinary copy once; $M(I,F)$ likewise counts copies, rather than
vertex embeddings.
For the empty edge set use $M(\varnothing)=M(I,\varnothing)=1$.
If $\boldsymbol I$ is a uniformly chosen copy of $I$ on $[n]$, and
$F_0$ is any fixed copy of an edge subset $F\subseteq I$, then
\begin{equation}\label{eq:copy-ratio}
 \PP(F_0\subseteq\boldsymbol I)=\frac{M(I,F)}{M(F)}.
\end{equation}
Indeed, permutation symmetry makes the probability independent of
the chosen copy $F_0$. Summing it over all $M(F)$ copies counts the
$M(I,F)$ copies contained in each $\boldsymbol I$.

\begin{lemma}[Graph hierarchy]\label{lem:hierarchy}
There is a disjoint probability tree on
$X=\binom{[n]}2$ whose path unions are copies of $H$ with its isolated
vertices removed, such that every level is $\sigma$-spread and its label
sizes $\ell_1,\ldots,\ell_k$ satisfy
\[
 1\le\ell_1,\qquad \ell_{i+1}\ge16\ell_i,
 \qquad \ell_k\le e(H).
\]
\end{lemma}

\begin{proof}
For a current nonempty edge set $I\subseteq H$, choose an edge
subset $S\subseteq I$ maximizing
\begin{equation}\label{eq:maximizer}
 R_\sigma(I,S)=\sigma^{-|S|}\PP(S_0\subseteq\boldsymbol I),
\end{equation}
where $S_0$ is any fixed copy of $S$ on $[n]$. The empty set is
allowed and has value $1$. Fix any deterministic rule for breaking
ties.

Write $h_I=|I|$ and $r=|S|>0$. The expectation constraint gives
$M(S)q^r\ge1/2$, while $M(I,S)\le\binom{h_I}{r}$.
By \eqref{eq:copy-ratio},
\[
 R_\sigma(I,S)\le2\binom{h_I}{r}(q/\sigma)^r
           =2\binom{h_I}{r}128^{-r}.
\]
For $r\ge h_I/17$, the bound
$\binom{h_I}{r}\le(\ee h_I/r)^r$ gives
\[
 R_\sigma(I,S)\le2(17\ee/128)^r<1.
\]
Thus every maximizing predecessor satisfies $|S|<|I|/17$.

Repeat the selection until the empty set is reached, and write
the resulting nested sequence in reverse order as
\[
 \varnothing=H_0\subset H_1\subset\cdots\subset H_k=H,
 \qquad |H_{i-1}|<|H_i|/17.
\]
Here and below $H$ denotes its edge set. Build a tree rooted at the
empty set. From each reached copy $S_0$ of $H_{i-1}$, take as
children all copies of $H_i$ containing $S_0$, with the uniform
distribution, and label the arc to a child $I_0$ by $I_0\setminus S_0$.
Such children exist, since every copy of $H_{i-1}$ extends to a copy
of $H_i$ on $[n]$. Keep distinct histories as distinct tree nodes.
The labels on every path are disjoint and their union is a copy of
$H$. Their sizes are
\[
 \ell_i=|H_i|-|H_{i-1}|>0,
 \qquad
 \ell_{i+1}>16|H_i|\ge16\ell_i.
\]

It remains to check conditional spread. At a node $S_0$, let $A$
be its random added label. If $J\cap S_0\ne\varnothing$, then
$\PP(J\subseteq A)=0$. Otherwise, for a uniform copy
$\boldsymbol I$ of $H_i$,
\begin{equation}\label{eq:conditional-ratio}
 \PP(J\subseteq A)=
 \frac{\PP(S_0\cup J\subseteq\boldsymbol I)}
      {\PP(S_0\subseteq\boldsymbol I)}.
\end{equation}
If the numerator is positive, $S_0\cup J$ is isomorphic to an edge
subset $U\subseteq H_i$. Maximality of $H_{i-1}$ in
\eqref{eq:maximizer} gives
\[
 \sigma^{-(|S_0|+|J|)}\PP(S_0\cup J\subseteq\boldsymbol I)
 \le \sigma^{-|S_0|}\PP(S_0\subseteq\boldsymbol I).
\]
Together with \eqref{eq:conditional-ratio}, this gives
$\PP(J\subseteq A)\le\sigma^{|J|}$. If the numerator is zero, the same
bound is immediate. This proves the lemma.
\end{proof}

\begin{proof}[Proof of Theorem~\ref{thm:main}]
Apply Theorem~\ref{thm:tree} to the tree of
Lemma~\ref{lem:hierarchy}, with $\sigma=128q$. Since
$\ell_k\le h=e(H)$, a Bernoulli edge set with parameter
\[
 r=\min\{1,\;2048\ee^{50}q(1+\log_2 h)\}
\]
contains a copy of $H$ with probability at least $2/3$. This proves
\eqref{eq:stronger}. Finally, $h\le\binom n2\le n^2$ and $n\ge2$
give $1+\log_2 h\le3\log_2 n$, and hence the stated constant
$C=6144\ee^{50}$.
\end{proof}

